\documentclass[10pt,a4paper]{amsart}
\usepackage[margin=19mm]{geometry}
\usepackage{amsmath,amssymb,mathtools}
\usepackage[scr=rsfs,cal=euler]{mathalfa}
\usepackage{xfrac}
\usepackage[unicode,hidelinks,bookmarks=false]{hyperref}
\newcommand{\ie}{i.e.\ }
\newcommand{\cf}{cf.\ }
\newcommand{\resp}{respectively\ }
\newcommand{\Subsection}{\subsection}
\providecommand{\nobreakdash}{\nobreak\hbox{-}}
\setlength{\emergencystretch}{2em}
\multlinegap0pt
\usepackage{amssymb}
\usepackage{tikz}
\usepackage[shortlabels]{enumitem}
\newtheorem{lem}{Lemma}
\newcommand{\vep}{\varepsilon}
\title{Regularity theory for fully nonlinear oblique transmission problems}
\author{I\~nigo U. Erneta}
\date{Source: arXiv:2606.14858v2 (CC BY 4.0)}
\begin{document}
\maketitle

\noindent\emph{Source and licence.} Regularity theory for fully nonlinear oblique transmission problems. Version arXiv:2606.14858v2 (CC BY 4.0), distributed by arXiv under the Creative Commons Attribution 4.0 International licence, http://creativecommons.org/licenses/by/4.0/, as recorded in the arXiv licence metadata for this exact version and shown on the abstract page https://arxiv.org/abs/2606.14858v2. Full licence notice: \texttt{licenses/arxiv-2606.14858-CC-BY-4.0.txt}.\par\smallskip

\noindent\textit{Editorial scope.} Counted unit: the article's lem lem (source0.tex lines 2739--2744). The article's complete original proof of it is delivered with it (source0.tex lines 2746--2847, one \texttt{proof} environment of 102 lines). No mathematics is written, repaired or re-derived: the statement and the proof are the article's own lines, in the article's own notation, and no retained line is substituted or re-flowed.  No further source-local context is needed: every cross-reference of every retained line resolves inside the two delivered spans.  Nothing was omitted from any retained span, so the package carries no omission record.  No retained line contains a citation, so the wrapper carries no bibliography and no bibliography entry.  No retained line contains a figure environment, an image-inclusion command or drawing code, so no image is delivered and the package declares no asset dependency.  Editorial formatting only, in addition: an AMS \texttt{amsart} preamble that restates the article's own declarations and macros used by the retained lines, namely the environments \texttt{lem} on separate counters, together with the standard packages those lines need.  The retained environments print in this document as Lemma 1 in order of appearance, whereas the article prints the same items with the numbering of its own full document.  The complete native source of the article as distributed in the arXiv e-print for this exact version is bundled unchanged as \texttt{sources/202811/source0.tex}.
\begin{lem}
There is a modulus of continuity $\widetilde{\omega}$ depending only on $\omega$, $\{\omega_{\kappa}\}_{\kappa > 0}$, and $\omega_0$ such that
\[
|u(x) - u(x_0)| \leq \widetilde{\omega}(|x-x_0|) \quad \text{ for all } x \in \overline{B}_{1} \text{ and } x_0 \in \partial B_1.
\]
\end{lem}

\begin{proof}
Let $\vep > 0$.
We must find a $\delta > 0$ depending only on $\vep$ and the given moduli, such that
\begin{equation}
\label{tri0}
|u(x) - u(x_0)| \leq \vep \quad \text{ whenever } |x - x_0| \leq \delta, \text{ with } x \in B_1 \text{ and } x_0 \in \partial B_1.
\end{equation}

\medskip

Assume that $|x - x_0| \leq \delta$, with $0 < \delta \leq 1/2$ small to be determined below.
By triangle inequality
\begin{equation}
\label{tri1}
|u(x) - u(x_0)| \leq |u(x) - u(\pi[x])| + |u(\pi[x]) - u(x_0)|.
\end{equation}

Since $|x| \geq |x_0| - |x- x_0| \geq 1/2$ and $d_x \leq |x- x_0|$, we have  that
\[
|\pi[x] - x_0 | = 
\left|\frac{x}{|x|} - x_0\right| 
= \frac{1}{|x|}\big|x - x_0 + (1-|x|) x_0\big|
\leq \frac{|x- x_0|}{|x|} +\frac{d_x}{|x|}
\leq 4 \delta
\]
and the second term on the right-hand side of~\eqref{tri1} can be controlled by the modulus on $\partial B_1$ as
\begin{equation}
\label{tri2}
|u(\pi[x]) - u(x_0)| \leq \omega(4 \delta).
\end{equation}

\medskip

To bound the first term, fix $0 < \kappa_0 \leq 1/4$ small to be chosen below. 
We distinguish two cases:
\begin{itemize}
\item If $|x_n| \geq \kappa_0$, then by the modulus $\omega_{\kappa_0}$, we have
\begin{equation}
\label{tri3}
|u(x) - u(\pi[x])| \leq \omega_{\kappa_0}(d_x) \leq \omega_{\kappa_0}(\delta).
\end{equation}
\item If $|x_n| \leq \kappa_0$, then we have the additional bounds
\begin{equation}
\label{eq:bod1}
|x'| \geq |x| - |x_n| \geq 1/2 - \kappa_0 \geq 1/4,
\end{equation}
and
\begin{equation}
\label{eq:bod2}
1 - |x'| \leq 1 - |x| + |x_n| < \delta + \kappa_0 \leq 2 \kappa_0.
\end{equation}
By triangle inequality, using the modulus $\omega_0$
\begin{equation}
\label{bodo}
\begin{split}
|u(x) - u(\pi[x])| &\leq \textstyle \left|u(x) - u(\frac{x'}{|x'|},0)\right| + \left| u(\pi[x]) - u(\frac{x'}{|x'|},0)\right|\\
& \quad \leq \textstyle 
\omega_0\left(\left|x -(\frac{x'}{|x'|},0) \right|\right) +  \omega_0\left(\left|\frac{x}{|x|} - (\frac{x'}{|x'|},0) \right| \right).
\end{split}
\end{equation}
Note that by~\eqref{eq:bod2}, we have
\begin{equation}
\label{cris1}
\textstyle
\left|x -(\frac{x'}{|x'|},0) \right| = \sqrt{(1- |x'|)^2 +x_n^2}
\leq \sqrt{5}\kappa_0 < 2 \sqrt{2} \kappa_0.
\end{equation}
On the other hand, by concavity of the square root
\[
|x| = \sqrt{|x'|^2+x_n^2} \leq |x'| + \frac{1}{2 |x'|}x_n^2,
\]
whence, by~\eqref{eq:bod1}, we deduce
\begin{equation}
\label{cris2}
\textstyle
\left|\frac{x}{|x|} - (\frac{x'}{|x'|},0) \right| = \sqrt{2}\sqrt{\frac{|x|-|x'|}{|x|}}
\leq \frac{|x_n|}{\sqrt{|x||x'|}}
\leq 2 \sqrt{2} \kappa_0.
\end{equation}
Applying~\eqref{cris1} and~\eqref{cris2} in~\eqref{bodo}, we conclude that
\begin{equation}
\label{tri4}
\begin{split}
|u(x) - u(\pi[x])| 
\leq 2 \omega_0(2 \sqrt{2} \kappa_0).
\end{split}
\end{equation}
\end{itemize}

\medskip

To finish the argument,
first choose $\kappa_0 > 0$ such that 
\[
2 \omega_0(2 \sqrt{2}\kappa_0) \leq \vep/2
\]
and then choose $\delta > 0$ such that
\[
\max\{ \omega(4\delta), \omega_{\kappa_0}(\delta) \} \leq \vep/2.
\]
Combining~\eqref{tri2},~\eqref{tri3},~and~\eqref{tri4} to control~\eqref{tri1} now yields the claim~\eqref{tri0}.
\end{proof}

\end{document}
